%
\chapter*{\centering{Abstract}}
\addcontentsline{toc}{chapter}{Abstract}

Gateways that relay Message Queuing Telemetry Transport traffic for the Internet of Things see attacks that few sites observe alone, yet their raw packets often cannot be shared to train an intrusion detection system. This thesis measures how far several gateways can jointly produce synthetic training data under formal privacy protection. A framework was built around a federated generator that releases noisy count tables over three rounds and samples packets from a dependency tree, with the noise split across gateways through secure aggregation. A federated conditional variational autoencoder trained with differentially private gradient descent served as the baseline. The study used a public denial-of-service traffic dataset of 39 million packets, split by capture, with five gateways holding skewed class shares, three random seeds, and a paired bootstrap decision rule. Detectors trained only on synthetic data reached a macro-averaged F1 score of 0.366, 0.425 and 0.433 at privacy budgets of 1, 5 and 10, against 0.448 for a random forest on pooled real data. At every budget the generator exceeded the autoencoder by 0.092 to 0.191, using at most 1.33 megabytes per run against 8.4 to 191 megabytes. When real data could be pooled, synthetic data gave no useful gain and simple class rebalancing performed better. A gateway lacking some classes gained about 0.04 from locally added synthetic rows. Protecting a whole connection cost 0.03 to 0.05, protecting a whole capture was not usable, and a budget of 5 grew to 14.2 with one honest gateway. A membership attack on the released tables showed a measurable benefit of the added noise. Generated data help most where pooling is impossible, and the choice of protection should follow the trust and bandwidth requirements of each deployment.

\textbf{Keywords:} intrusion detection, Internet of Things, federated learning, differential privacy, secure aggregation, synthetic data generation, MQTT
